\documentclass[a4paper,12pt]{article} % Tipo de documento y tamaño de letra.
\usepackage[utf8]{inputenc} % Para que reconozca caracteres especiales como tildes.
\usepackage[spanish, es-nodecimaldot]{babel} % Para definir el idioma y el formato de fechas en español.
\usepackage{amsmath, amssymb} % Paquetes para símbolos y ecuaciones avanzadas (Para hacer una ecuación con un recuadro).
\usepackage{ragged2e} % Para justificar textos.
\usepackage{multicol} % Para crear columnas.
\usepackage{slashed} % Para generar números o variables canceladas (Una línea diagonal sobre el número o variable).
\usepackage{graphicx} % Paquete para poner graficas.
\usepackage{float} % Para poder dejar la imagen en un punto fijo del documento.
\usepackage{longtable} % Para tablas largas.
\usepackage{cancel} % Para hacer cancelaciones en medio de operaciones.
\usepackage{xcolor} % Para el color y esas cosas.
\usepackage{apacite} % Paquete para citas APA.
\usepackage{fancyhdr} % Para personalizar encabezados y pies de página.
\usepackage{geometry} % Para ajustar márgenes.
\usepackage{setspace} % Para ajustar el interlineado.
\usepackage{natbib} % Para manejar las citas y que no queden pegadas y feas :P.
\usepackage{hyperref} % Para que los enlaces funcionen con su URL.
\usepackage[version=4]{mhchem}  % Versión recomendada para mayor funcionalidad.
\usepackage{chemfig} %También para hacer estructuras químicas.
\usepackage{array} % Para mejor control de las columnas.
\usepackage{siunitx} % Para unidades físicas.


\pagestyle{fancy}
\fancyhf{} % Limpia encabezados y pies de página predeterminados
\fancyhead[L]{\textbf{LABORATORIO QUICKFIELD.}} % Encabezado izquierdo
\fancyhead[R]{\textbf{ELECTROMAGNETISMO I.}} % Encabezado derecho
\fancyfoot[C]{\thepage} % Número de página en el centro del pie

\geometry{headheight=12.96115pt, top=1in, bottom=1in, left=1in, right=1in}
\begin{document}


\begin{figure}[H]
	\centering
	\includegraphics[width=0.5\textwidth]{Unir.jpg}
\end{figure}

\vspace{\baselineskip}

% Ajustamos manualmente todo ;)
\begin{center}
	\textbf{\Large } \\[0.9cm]
	\vspace{\baselineskip}
	\textbf{Liz Ángel Núñez Torres.} \\[0.6cm]
	\textbf{Grado en Física.} \\[0.6cm]
	\textbf{Universidad Internacional de La Rioja.} \\[0.6cm]
	\textbf{8 de Diciembre 2025.}
\end{center}


\newpage
\tableofcontents

\newpage

\section{Hilo Infinito.}

\begin{justify}
Hacemos el \textbf{Hilo Infinito} con el mallado y la solución del mismo.
\end{justify}

\begin{justify}
    Configuración en \textbf{QuickField}:
\end{justify}

\begin{multicols}{3}
    \begin{itemize}
        \item \( \lambda = 2 \times 10^{-9}. \)
        
        \item \( \epsilon_r =1 .\)
        
        \item \( \mathrm{V} = 0. \)
    \end{itemize}
\end{multicols}


\begin{figure}[H]
    \centering
    \includegraphics[width=0.8 \textwidth]{figura1.jpg}
\end{figure}

\begin{justify}
    Solucionamos y obtenemos los resultados.
\end{justify}

\begin{figure}[H]
    \centering
    \includegraphics[width=1.1 \textwidth]{figura1.3.jpg}
\end{figure}

\newpage

\begin{justify}
    Modificamos el script en \textbf{Python} para nuestro ejemplo.
\end{justify}

\begin{figure}[H]
    \centering
    \includegraphics[width=0.9 \textwidth]{figura1.1.jpg}
\end{figure}

\begin{justify}
    Obtenemos el plot de las figuras.
\end{justify}

\begin{figure}[H]
    \centering
    \includegraphics[width=1.1 \textwidth]{figura1.2.jpg}
\end{figure}

\subsection{Conclusión.}

\begin{justify}
    Lo resultados en \textbf{QuickField} validan la dependencia radial teórica.
\end{justify}

\newpage

\section{Plano Cargado.}


\begin{justify}
    Configuración en \textbf{QuickField}:
\end{justify}

\begin{multicols}{3}
    \begin{itemize}
        \item Vacío \(\epsilon_r = 1.  \)
        
        \item Contorno \( \mathrm{V} = 0.\)
        
        \item Carga  \( \sigma = 0.5 \, C/m^2. \)
    \end{itemize}
\end{multicols}


\begin{figure}[H]
    \centering
    \includegraphics[width=0.8 \textwidth]{figura2.jpg}
\end{figure}


\begin{justify}
    Solucionamos el plano.
\end{justify}

\begin{figure}[H]
    \centering
    \includegraphics[width=1.1 \textwidth]{figura2.1.jpg}
\end{figure}

\newpage

\begin{justify}
    Corregimos \(\sigma\) en el script de Python.
\end{justify}

\begin{figure}[H]
    \centering
    \includegraphics[width=1 \textwidth]{figura2.2.jpg}
\end{figure}

\begin{justify}
    La graficas que obtenemos son:
\end{justify}

\begin{figure}[H]
    \centering
    \includegraphics[width=0.8 \textwidth]{figura2.3.jpg}
\end{figure}

\subsection{Conclusión.}

\begin{justify}
    Los resultados en QuickField validan que \textbf{el campo eléctrico de un plano infinito es constante e independiente de la distancia}, cumpliendo \(E = \sigma/(2 \epsilon_0)\).
\end{justify}

\newpage 

\section{Condensador Plano.}

\begin{justify}
    Configuración en \textbf{QuickField}:
\end{justify}


\begin{justify}
    \begin{itemize}
        \item \textbf{Dimensiones Condensador:}
    \end{itemize}
\end{justify}

\begin{multicols}{3}
    \begin{itemize}
        \item \( W =20 mm.\)
        
        \item \( h = 20 \, mm. \)
        
        \item \( L_z = 10 \, mm. \)
    \end{itemize}
\end{multicols}

\begin{multicols}{2}
    \begin{itemize}
        \item Material IZQ. \( \epsilon_{r1} = 4. \)
        
        \item  Material DER. \( \epsilon_{r2} = 1. \)
    \end{itemize}
\end{multicols}

\begin{multicols}{2}
    \begin{itemize}
        \item Voltaje \( V = 10\,V.  \)
        
        \item Separación \( d = 20\,mm.\)
    \end{itemize}
\end{multicols}


\begin{justify}
    \textbf{Condensador Plano y \(Q_s\)}
\end{justify}


\begin{figure}[H]
    \centering
    \includegraphics[width=0.8 \textwidth]{figura3.jpg}
\end{figure}

\begin{justify}
    Tabla QuickField.
\end{justify}

\begin{figure}[H]
    \centering
    \includegraphics[width=1.1 \textwidth]{figura3.1.jpg}
\end{figure}

\newpage

\begin{justify}
    Comprobamos \(\mathbf{Q_s}\) con Python.
\end{justify}

\begin{figure}[H]
    \centering
    \includegraphics[width=1 \textwidth]{figura3.2.jpg}
\end{figure}

\begin{justify}
    Gráficas del \textbf{desplazamiento} y el \textbf{campo eléctrico.}
\end{justify}

\begin{figure}[H]
    \centering
    \includegraphics[width=1 \textwidth]{figura3.3.jpg}
\end{figure}

\subsection{Conclusión.}

\begin{justify}
    Los resultados obtenidos mediante QuickField muestran un valor prácticamente constante de \(D_y\), en concordancia con la teoría, lo que confirma la correcta modelización del sistema.
\end{justify}

\vspace{\baselineskip}

\section{Espira Corriente.}

\begin{justify}
    Configuración en \textbf{QuickField}:
\end{justify}

\begin{multicols}{3}
    \begin{itemize}
        \item \( w = 80\)
        
        \item \( h = 60 \)
        
        \item \( y = 30 \)
    \end{itemize}
\end{multicols}


\begin{multicols}{3}
    \begin{itemize}
        \item Radio \( R = 12.5\)
        \item Corriente \( I = 2.5 \)
    \end{itemize}
\end{multicols}


\begin{figure}[H]
    \centering
    \includegraphics[width=0.9 \textwidth]{figura4.jpg}
\end{figure}

\begin{justify}
    Obtenemos la tabla de datos.
\end{justify}


\begin{figure}[H]
    \centering
    \includegraphics[width=0.9 \textwidth]{figura4.1.jpg}
\end{figure}


\newpage


\begin{justify}
    Modificamos los valores en Python.
\end{justify}

\begin{figure}[H]
    \centering
    \includegraphics[width=0.9 \textwidth]{figura4.2.jpg}
\end{figure}

\begin{justify}
    Obtenemos el Plot del código.
\end{justify}

\begin{figure}[H]
    \centering
    \includegraphics[width=0.9 \textwidth]{figura4.3.jpg}
\end{figure}

\subsection{Conclusión.}

\begin{justify}
Los resultados obtenidos mediante QuickField muestran una concordancia con la solución teórica. Se comprobo con \(Bz_{teo}\) en Python, que la componente axial del campo magnético sigue la depencia teórica.
\end{justify}

\section{Revisión.}

\begin{figure}[H]
    \centering
    \includegraphics[width=1 \textwidth]{revision.jpg}
\end{figure}

\end{document}